\documentclass[10pt,a4paper]{amsart}
\usepackage{amsmath,amssymb,amsthm,graphicx}
\usepackage{verbatim}
\newenvironment{reference}{\comment}{\endcomment}
%\newenvironment{reference}{}{}
\newenvironment{slogan}{\comment}{\endcomment}
\newenvironment{history}{\comment}{\endcomment}
% For commutative diagrams we use Xy-pic
\usepackage[all]{xy}

% We use 2cell for 2-commutative diagrams.
\xyoption{2cell}
% Theorem environments.
%
\theoremstyle{plain}
\newtheorem{theorem}[subsection]{Theorem}
\newtheorem{proposition}[subsection]{Proposition}
\newtheorem{lemma}[subsection]{Lemma}

\theoremstyle{definition}
\newtheorem{definition}[subsection]{Definition}
\newtheorem{example}[subsection]{Example}
\newtheorem{exercise}[subsection]{Exercise}
\newtheorem{situation}[subsection]{Situation}

\theoremstyle{remark}
\newtheorem{remark}[subsection]{Remark}
\newtheorem{remarks}[subsection]{Remarks}

\numberwithin{equation}{subsection}

% Macros
%
\def\lim{\mathop{\mathrm{lim}}\nolimits}
\def\colim{\mathop{\mathrm{colim}}\nolimits}
\def\Spec{\mathop{\mathrm{Spec}}}
\def\Hom{\mathop{\mathrm{Hom}}\nolimits}
\def\Ext{\mathop{\mathrm{Ext}}\nolimits}
\def\SheafHom{\mathop{\mathcal{H}\!\mathit{om}}\nolimits}
\def\SheafExt{\mathop{\mathcal{E}\!\mathit{xt}}\nolimits}
\def\Sch{\mathit{Sch}}
\def\Mor{\mathop{\mathrm{Mor}}\nolimits}
\def\Ob{\mathop{\mathrm{Ob}}\nolimits}
\def\Sh{\mathop{\mathit{Sh}}\nolimits}
\def\NL{\mathop{N\!L}\nolimits}
\def\CH{\mathop{\mathrm{CH}}\nolimits}
\def\proetale{{pro\text{-}\acute{e}tale}}
\def\etale{{\acute{e}tale}}
\def\QCoh{\mathit{QCoh}}
\def\Ker{\mathop{\mathrm{Ker}}}
\def\Im{\mathop{\mathrm{Im}}}
\def\Coker{\mathop{\mathrm{Coker}}}
\def\Coim{\mathop{\mathrm{Coim}}}

% Boxtimes
%
\DeclareMathSymbol{\boxtimes}{\mathbin}{AMSa}{"02}

%
% Macros for moduli stacks/spaces
%
\def\QCohstack{\mathcal{QC}\!\mathit{oh}}
\def\Cohstack{\mathcal{C}\!\mathit{oh}}
\def\Spacesstack{\mathcal{S}\!\mathit{paces}}
\def\Quotfunctor{\mathrm{Quot}}
\def\Hilbfunctor{\mathrm{Hilb}}
\def\Curvesstack{\mathcal{C}\!\mathit{urves}}
\def\Polarizedstack{\mathcal{P}\!\mathit{olarized}}
\def\Complexesstack{\mathcal{C}\!\mathit{omplexes}}
% \Pic is the operator that assigns to X its picard group, usage \Pic(X)
% \Picardstack_{X/B} denotes the Picard stack of X over B
% \Picardfunctor_{X/B} denotes the Picard functor of X over B
\def\Pic{\mathop{\mathrm{Pic}}\nolimits}
\def\Picardstack{\mathcal{P}\!\mathit{ic}}
\def\Picardfunctor{\mathrm{Pic}}
\def\Deformationcategory{\mathcal{D}\!\mathit{ef}}
\usepackage{cleveref}
\crefname{originallemma}{original lemma}{original lemmas}
\Crefname{originallemma}{Original lemma}{Original lemmas}
\makeatletter
\newlabel{cohomology-lemma-projection-formula}{{01E8}{source}{}{sourcepointer-1}{}}
\newlabel{cohomology-lemma-projection-formula@cref}{{[originallemma][0][]01E8}{[1][0][]source}{}{}{}}
\makeatother
\crefname{originallemma}{original lemma}{original lemmas}
\Crefname{originallemma}{Original lemma}{Original lemmas}
\makeatletter
\newlabel{lemma-c1-cap-additive}{{02SP}{source}{}{sourcepointer-2}{}}
\newlabel{lemma-c1-cap-additive@cref}{{[originallemma][0][]02SP}{[1][0][]source}{}{}{}}
\makeatother
\crefname{originallemma}{original lemma}{original lemmas}
\Crefname{originallemma}{Original lemma}{Original lemmas}
\makeatletter
\newlabel{lemma-flat-pullback-cap-c1}{{02SS}{source}{}{sourcepointer-3}{}}
\newlabel{lemma-flat-pullback-cap-c1@cref}{{[originallemma][0][]02SS}{[1][0][]source}{}{}{}}
\makeatother
\crefname{originallemma}{original lemma}{original lemmas}
\Crefname{originallemma}{Original lemma}{Original lemmas}
\makeatletter
\newlabel{lemma-gysin-back}{{0F95}{source}{}{sourcepointer-4}{}}
\newlabel{lemma-gysin-back@cref}{{[originallemma][0][]0F95}{[1][0][]source}{}{}{}}
\makeatother
\crefname{originallemma}{original lemma}{original lemmas}
\Crefname{originallemma}{Original lemma}{Original lemmas}
\makeatletter
\newlabel{lemma-linebundle}{{0B74}{source}{}{sourcepointer-5}{}}
\newlabel{lemma-linebundle@cref}{{[originallemma][0][]0B74}{[1][0][]source}{}{}{}}
\makeatother
\crefname{originallemma}{original lemma}{original lemmas}
\Crefname{originallemma}{Original lemma}{Original lemmas}
\makeatletter
\newlabel{lemma-relative-effective-cartier}{{02TR}{source}{}{sourcepointer-6}{}}
\newlabel{lemma-relative-effective-cartier@cref}{{[originallemma][0][]02TR}{[1][0][]source}{}{}{}}
\makeatother
\title{The zero section of a line bundle: regularity, pushforward, pullback and Chern class formulae}
\author{The Stacks Project Authors}
\date{stacks-project:chow.tex:local-snapshot-20261012; GNU FDL 1.2 or later}
\begin{document}
\maketitle
\noindent\textit{Editorial note.} All delivered mathematical text is the author's own, copied from the pinned Stacks Project chapter file chow.tex; no statement or proof step was rewritten, reordered or removed.\par\medskip
\noindent\textit{Editorial note.} Preamble provenance: the pinned chapter file opens with \textbackslash{}input\{preamble\} and carries no preamble text of its own, so the project's single Stacks-wide preamble staged at sources/stacks-chow/preamble.tex is excerpted verbatim as three spans, lines 8-12 (the verbatim package together with the project's own reference, slogan and history display-environment definitions, including the author's commented-out older form of the reference environment), lines 14-18 (Xy-pic setup) and lines 154-217 (the native theorem declarations and mathematical macros). The lines 8-12 span is required because a counted statement or its surrounding text may use those display environments, which the project typesets as comments. Nothing was added, removed or redefined; the wrapper class and the amsmath/amssymb/amsthm/graphicx base packages are supplied by the assembler.\par\medskip
\noindent\textit{Editorial note.} Same-chapter pointers are located at their \textbackslash{}label lines inside the pinned chapter file; cross-chapter pointers are located in the project's own tag list staged in the snapshot (sources/stacks-chow/tags.txt), at the verbatim lines carrying their canonical tags.\par\medskip
\noindent\textit{Editorial note.} The author's sentence at source lines 5801-5802, 'We omit the proof that the vanishing locus of this section is precisely o(X).', is retained verbatim inside the counted proof; the identification of the section and the use of the projection formula are inside the span, and only that routine check is left to the reader.\par\medskip
\noindent\textit{Allowed standard prerequisites (editorial).}\begin{itemize}
\item The projection formula and the structure of the pushforward of the structure sheaf of a line bundle. Cohomology, label cohomology-lemma-projection-formula, invoked at proof line 5797. Identifies the pushforward of the inverse of the pulled back invertible sheaf with the symmetric algebra twisted by the inverse, which is what makes the canonical trivialising section visible.
\item The Gysin pullback of an effective Cartier divisor and additivity and flat pullback compatibility of capping with a first Chern class. Same chapter chow.tex, labels lemma-relative-effective-cartier (proof lines 5807 and 5811), lemma-c1-cap-additive (proof line 5812) and lemma-flat-pullback-cap-c1 (proof line 5814). The first two formulae of that lemma give the identity o\textasciicircum{}* p\textasciicircum{}* = id and the expression of the pushforward of a class as minus the pullback of the cap with the Chern class, and the last two let the Chern class be moved through the flat pullback to derive the remaining identities.
\item The computation of the composite of the pushforward with the pullback along the zero section. Same chapter chow.tex, label lemma-gysin-back, invoked at proof line 5817. Combined with the identity o\textasciicircum{}*p\textasciicircum{}*L = L, it gives that the composite of the pullback with the pushforward is minus the cap with the Chern class.
\end{itemize}\medskip
\ifdefined\hypertarget\hypertarget{sourcepointer-1}{}\fi
\noindent\textit{Source reference.} Original lemma 01E8: sources/stacks-chow/tags.txt, source lines 1825--1825. Navigation pointer for \textbackslash{}ref\{cohomology-lemma-projection-formula\} used inside the counted span (line 5797); the argument that consumes it is retained in the span.\par\medskip
\ifdefined\hypertarget\hypertarget{sourcepointer-2}{}\fi
\noindent\textit{Source reference.} Original lemma 02SP: sources/stacks-chow/pinned-original-source.tex, source lines 4330--4330. Navigation pointer for \textbackslash{}ref\{lemma-c1-cap-additive\} used inside the counted span (line 5812); the argument that consumes it is retained in the span.\par\medskip
\ifdefined\hypertarget\hypertarget{sourcepointer-3}{}\fi
\noindent\textit{Source reference.} Original lemma 02SS: sources/stacks-chow/pinned-original-source.tex, source lines 4533--4533. Navigation pointer for \textbackslash{}ref\{lemma-flat-pullback-cap-c1\} used inside the counted span (line 5814); the argument that consumes it is retained in the span.\par\medskip
\ifdefined\hypertarget\hypertarget{sourcepointer-4}{}\fi
\noindent\textit{Source reference.} Original lemma 0F95: sources/stacks-chow/pinned-original-source.tex, source lines 5463--5463. Navigation pointer for \textbackslash{}ref\{lemma-gysin-back\} used inside the counted span (line 5817); the argument that consumes it is retained in the span.\par\medskip
\ifdefined\hypertarget\hypertarget{sourcepointer-5}{}\fi
\noindent\textit{Source reference.} Original lemma 0B74: sources/stacks-chow/pinned-original-source.tex, source lines 5737--5737. Navigation pointer for \textbackslash{}ref\{lemma-linebundle\} used inside the counted span (line 5776); the argument that consumes it is retained in the span.\par\medskip
\ifdefined\hypertarget\hypertarget{sourcepointer-6}{}\fi
\noindent\textit{Source reference.} Original lemma 02TR: sources/stacks-chow/pinned-original-source.tex, source lines 5598--5598. Navigation pointer for \textbackslash{}ref\{lemma-relative-effective-cartier\} used inside the counted span (line 5807); the argument that consumes it is retained in the span.\par\medskip
\begin{lemma}
\label{lemma-linebundle-formulae}
In the situation of Lemma \ref{lemma-linebundle} denote $o : X \to L$
the zero section (see proof of the lemma). Then we have
\begin{enumerate}
\item $o(X)$ is the zero scheme of a regular global section of
$p^*\mathcal{L}^{\otimes -1}$,
\item $o_* : \CH_k(X) \to \CH_k(L)$ as $o$ is a closed immersion,
\item $o^* : \CH_{k + 1}(L) \to \CH_k(X)$ as $o(X)$
is an effective Cartier divisor,
\item $o^* p^* : \CH_k(X) \to \CH_k(X)$ is the identity map,
\item $o_*\alpha = - p^*(c_1(\mathcal{L}) \cap \alpha)$ for any
$\alpha \in \CH_k(X)$, and
\item $o^* o_* : \CH_k(X) \to \CH_{k - 1}(X)$ is equal to the map
$\alpha \mapsto - c_1(\mathcal{L}) \cap \alpha$.
\end{enumerate}
\end{lemma}

\begin{proof}
Since $p_*\mathcal{O}_L = \text{Sym}^*(\mathcal{L})$ we have
$p_*(p^*\mathcal{L}^{\otimes -1}) =
\text{Sym}^*(\mathcal{L}) \otimes_{\mathcal{O}_X} \mathcal{L}^{\otimes -1}$
by the projection formula
(Cohomology, Lemma \ref{cohomology-lemma-projection-formula})
and the section mentioned in (1) is
the canonical trivialization $\mathcal{O}_X \to
\mathcal{L} \otimes_{\mathcal{O}_X} \mathcal{L}^{\otimes -1}$.
We omit the proof that the vanishing locus of
this section is precisely $o(X)$. This proves (1).

\medskip\noindent
Parts (2), (3), and (4) we've seen in the course of the proof of
Lemma \ref{lemma-linebundle}. Of course (4) is the first
formula in Lemma \ref{lemma-relative-effective-cartier}.

\medskip\noindent
Part (5) follows from the second formula in
Lemma \ref{lemma-relative-effective-cartier},
additivity of capping with $c_1$ (Lemma \ref{lemma-c1-cap-additive}),
and the fact that capping with $c_1$ commutes with flat pullback
(Lemma \ref{lemma-flat-pullback-cap-c1}).

\medskip\noindent
Part (6) follows from Lemma \ref{lemma-gysin-back}
and the fact that $o^*p^*\mathcal{L} = \mathcal{L}$.
\end{proof}

\end{document}
